\section{Introduction}
\label{sec:intro}

The cloud computing industry is built on a fundamental inversion of ownership.
When a developer deploys an application to AWS, GCP, or Azure, the provider
assumes custodial control over the application's data, encryption keys, network
identity, and operational continuity. The user pays for access to their own data
through the provider's proprietary APIs. Switching providers requires rewriting
application code, migrating data through provider-controlled export tools (when
they exist), and accepting downtime measured in days or weeks. This is not a
service market. It is a custody arrangement.

The consequences are well-documented. Amazon Web Services holds approximately 31\%
of global cloud infrastructure spend \cite{synergy2024cloud}. Google Cloud and
Microsoft Azure hold a combined 33\%. Together, three companies control the
operational infrastructure for the majority of the internet's applications. When
AWS experiences an outage, large fractions of the internet become unavailable
\cite{aws2021outage}. When a provider decides to change pricing, deprecate a
service, or restrict access, customers have no recourse except migration---which
the provider has made as expensive as possible.

This paper presents an alternative: the Lux provider model, in which cloud
infrastructure providers are service markets, not data owners. The key insight
is a separation of concerns:

\begin{enumerate}[label=(\roman*)]
    \item \textbf{Users own their keys.} All data is encrypted client-side with
    keys stored in capsule vaults. No provider ever holds plaintext data or
    decryption keys.
    \item \textbf{Providers compete on service quality.} Relay latency, storage
    durability, compute throughput, and recovery speed are the competitive
    dimensions---not data lock-in.
    \item \textbf{Exit is a first-class right.} Users can switch providers at
    any time without data loss, because providers never hold the only copy of
    user data in a format only they can read.
    \item \textbf{SLAs are chain-enforced.} Service level agreements are encoded
    as smart contracts with stake-backed penalties for violations, verified by
    on-chain receipts.
\end{enumerate}

The remainder of this paper formalizes this model. Section~\ref{sec:model}
defines the five provider roles. Section~\ref{sec:reputation} introduces
the reputation protocol. Section~\ref{sec:sla} describes SLA enforcement.
Section~\ref{sec:markets} defines the compute and resource markets.
Section~\ref{sec:portability} proves the portability guarantee.
Section~\ref{sec:comparison} compares against centralized clouds.
Section~\ref{sec:economics} presents the economic model.
Section~\ref{sec:related} discusses related work.
Section~\ref{sec:conclusion} concludes.

